% !PRINTED 247
\weberpage{247}
Richtet man die Funktionen $\Phi_{a,c,\partial}$, etwa durch
geeignete Bestimmung von Konstanten, die darin noch verfügbar sind,
so ein, daß sie auch für ein unendliches imaginäres $\omega$, d.\,h.
für $q=0$ endlich bleiben, so sind diese Funktionen auch für ein
unendliches $j(\omega)$ endlich, und die Koeffizienten in (6) sind
konstant. Dies ist aber nur dadurch möglich, daß die
$\Phi_{a,c,\partial}$ \gesperrt{alle einer und derselben Konstanten $C$ gleich
sind}. Kennt man $\Phi_{a,c,\partial}$ als rationale Funktion einer
anderen Größe $\Psi_{a,c,\partial}$, so ist $\Phi_{a,c,\partial} = C$
entweder eine \gesperrt{Identität} oder eine \gesperrt{Transformationsgleichung} für
$\Psi_{a,c,\partial}$.

\begin{center}
\textbf{§ 71. \gesperrt{Die Transformationsgleichungen für $\gamma_2$ und $\gamma_3$.}}
\end{center}

Transformationsgleichungen erster Stufe erhält man zunächst aus der
Betrachtung der Funktionen § 54, (4), (5):
\[
\gamma_2(\omega) = \sqrt[3]{j(\omega)}, \quad \gamma_3(\omega) = \sqrt{j(\omega)-1728}.
\]
Wenn $n$ nicht durch 3 teilbar ist, so können wir die Zahlen $a$,
$c$, $\partial$ so wählen, daß \gesperrt{immer $c$ durch 3 teilbar
wird}.

Nun übt nach § 54, (15) eine lineare Transformation
$\begin{pmatrix} \alpha, & \beta \\ \gamma, & \delta \end{pmatrix}$
auf die Funktion $\gamma_2(\omega)$ den Einfluß:
\[
\gamma_2\left(\frac{\gamma+\delta\omega}{\alpha+\beta\omega}\right) = e^{-\frac{2\pi i}{3}(\alpha\gamma+\beta\delta-\alpha\beta-\alpha^2\beta\delta)} \gamma_2(\omega).
\]
In den Zusammensetzungen (8), (10) des § 69 wird
\[
\lambda \equiv n, \quad \delta \equiv 0 \ (\mathrm{mod}\ 3),
\]
und dann kann man $\alpha$ noch so bestimmen, daß auch $\alpha$ und
folglich $c_2$ durch 3 teilbar werden.

Daraus ergibt sich, daß durch die beiden Substitutionen
\[
(\omega, \omega+1), \quad \left(\omega, -\frac{1}{\omega}\right)
\]
die $\nu$ Funktionen
\weq{1}{\gamma_2\left(\frac{c+\partial\omega}{a}\right) \gamma_2(\omega)^{-n}}
nur untereinander vertauscht werden und also die Wurzeln einer
Transformationsgleichung erster Stufe sind.

Ist zweitens $n$ eine \gesperrt{ungerade} Zahl, so nehme man $c$
\gesperrt{gerade} an.
